\section{Protocol}
\label{sec:protocol}

\textbf{Data and split.} HAM10000~\cite{tschandl2018ham10000} contains 10,015 dermoscopic images of 7,470
lesions in seven diagnostic classes, each with a binary lesion mask. We use a lesion-grouped, stratified
80/10/10 split with a fixed split seed, so that every training seed is evaluated on the identical
\todo{1,002}-image test set (image ids checked identical across runs). Images are resized to
$224\times224$ and normalized with each backbone's ImageNet statistics.

\textbf{Models and training.} Table~\ref{tab:models} lists the two backbones. They differ in capacity and
pretraining data, so we treat architecture differences as differences between these specific pretrained
models, not as effects of inductive bias. Each model is fine-tuned in two phases, head only and then the
full network~\cite{kumar2022finetuning}, with AdamW~\cite{loshchilov2019adamw} and class-weighted cross
entropy. Hyperparameters were chosen on validation macro-F1 (search grid in Appendix \TODO{A, if it fits the 9-page limit}, otherwise in the code release); the checkpoint
of each run is the epoch with the best validation macro-F1, and the test split is evaluated once per
selected checkpoint. We train three seeds (42, 43, 44) per model.
\TODO{T2: state ViT phase-1 lr 5e-2 and extend the grid description once the seed-42 reproduction check passes.}

\begin{table}[t]
\centering\small
\caption{Backbones and selected configurations (from configs/*.yaml). \TODO{verify parameter counts and
pretraining data from the model cards and log them in PROVENANCE.md.}}
\label{tab:models}
\begin{tabular}{lcc}
\toprule
 & EfficientNet-B0 & ViT-B/16 \\
\midrule
Parameters & \todo{5.3M} & \todo{86M} \\
Pretraining data & \todo{ImageNet-1k} & \todo{ImageNet-21k, fine-tuned 1k} \\
Phase 1 / 2 lr & $10^{-3}$ / $8\times10^{-4}$ & \todo{$5\times10^{-2}$} / $4\times10^{-5}$ \\
Batch, weight decay & 16, 0.01 & 8, 0.1 \\
Dropout, augmentation & 0.2 (+ drop-connect 0.4), yes & 0.2, no \\
\bottomrule
\end{tabular}

\end{table}

\textbf{Explainers.} Grad-CAM~\cite{selvaraju2017gradcam} on the last activation of EfficientNet-B0 and
attention rollout~\cite{abnar2020attention} (mean head fusion, residual weight 0.5, no discarding) on
ViT-B/16 are the customary explainers. Rollout is class-agnostic. As a shared, class-specific explainer we
use Integrated Gradients~\cite{sundararajan2017axiomatic} on both models: 32 steps, an all-zero baseline in
normalized input space (the per-channel mean color), the predicted class as target, and the absolute
attribution summed over color channels. All heatmaps are $224\times224$ and min-max normalized per image.

\textbf{Controls.} \emph{Random}: an i.i.d.\ uniform heatmap seeded by the image id, identical across models
and seeds. \emph{Center prior}: a fixed heatmap that decreases with squared distance from the image center
and is identical for every image; it uses no model and no image content.

\textbf{Faithfulness.} Following RISE~\cite{petsiuk2018rise}, deletion replaces pixels, in heatmap order
and 512 pixels per step, with a Gaussian-blurred copy (kernel 11, $\sigma=5$), and insertion starts from a
mean-color canvas and restores the original pixels. We record the probability of the class predicted on
the clean image and report the area under each curve, macro-averaged over the seven predicted classes.
Lower deletion and higher insertion area indicate a more faithful ordering. The same target class and
images are used for every ordering.

\textbf{Localization.} Masks are resized to $224\times224$ with nearest-neighbor interpolation. IoU uses the
top 20\% of heatmap pixels. The pointing game counts a hit when the heatmap maximum lies inside the mask;
we report the exact-pixel version as primary and a 15-pixel tolerance~\cite{zhang2016topdown} as
secondary.

\textbf{Statistics.} For each run-level metric we report the mean $\pm$ SD over the three seeds. For
differences between two conditions we average each image's value over seeds and compute a paired
bootstrap 95\% confidence interval over images (10,000 resamples).
